\section{Experiments}

% NOTE: this chapter describes the configuration that produced
% results/final_ns3_grid.json (378 ns-3 simulations, 0 failures, dirs results/ns3f_*).
% Launchers: train_matched_hparams.sh (training), run_ns3_final.sh (evaluation).

This section outlines the experimental setup based on the proposed methodology. It details the training and evaluation procedures for the Multi-Agent Reinforcement Learning computational environment. The resources used to represent the backbone network, the baseline routing methods, and the metrics employed to assess network performance are described. Additionally, the parameters utilized for model training are presented. The aim of this section is to ensure that the methodology can be reproduced in other environment.

\subsection{Computing Environment}
\label{sec:setup-compute}

The experiment was conducted on a single Linux server (Rocky 9) equipped with an AMD EPYC 7543 processor (128 hardware threads) and 512 GB of RAM. We use the CPU to train and evaluate the model with concurrent core parallelization. The software stack included Python 3.9, ns-3.42, PyTorch 2.6.0, PyTorch Geometric 2.6.1, NetworkX, and Stable-Baselines3 2.7.1. 

\subsection{Networks and Traffic Resources}
\label{sec:setup-instances}

The simulation utilizes a set of topologies and traffic patterns. We use seventeen topologies from SNDLib to train the model, except Abilene, GEANT, and Germany50. The training set includes topologies ranging from 10 to 54 nodes and 21 to 80 edges. This approach aims to enable the model’s ability to generalize across different network structures. In contrast, traffic data require a different approach because measured traffic is only available for Abilene, GEANT, and Germany50.

To overcome the limitations of SNDLib resources for policy training, we generate synthetic traffic using the modulated-gravity model with the TMgen library \citep{internet-traffic}. This model uses an analogy with Newton’s law of gravitation to represent traffic between nodes and is the standard tool for this purpose. For each topology,  traffic patterns are generated and evaluated at five volume scales: $\{0.6, 0.8, 1.0, 1.2, 1.5\}$. We normalize these so that a scale of $1.0$ sets the OSPF bottleneck at exactly $100\%$. This methodology allows us to evaluate both feasible and congested traffic conditions.

Additionally, the measured traffic in SNDLib is insufficient to utilize the capacity of the backbone topologies fully. To simulate scenarios where link capacities are exceeded, demand scaling variations are applied to the evaluation topologies.  The three backbones span $12$, $22$ and $50$ nodes. Full instance properties and the complete evaluation protocol are given in Table~\ref{tab:eval-setup} of Appendix~\ref{app:instances}.


\subsection{Methods Compared}
\label{sec:setup-methods}

MARL is compared against two deployed routing protocols and a single-agent reinforcement learning controller, all evaluated on identical inputs.

\begin{enumerate}
    \item \textbf{OSPF.} The OSPF cost metric is defined as $\text{cost}(e) = \text{refBW}/\text{cap}(e)$. This configuration is commonly deployed in practice, particularly in networks with heterogeneous link capacities.
    \item \textbf{ECMP.} Equal-Cost Multipath over the paths that the same cost metric leaves tied, assigned one path per flow by hashing the packet header rather than divided fractionally. Because the tie is on cost, ECMP inherits OSPF's capacity awareness and forms a stronger baseline than OSPF alone.
    \item \textbf{Single-agent GNN.} A centralized controller that observes the entire network and selects one of $k=3$ candidate paths per demand. The network state is represented using a graph neural network (GNN) with shared weights across all links.
    \item \textbf{MARL.} Multi-Agent Reinforcement Learning (MARL) assigns an agent to each network node. Paths are constructed incrementally, with each agent making decisions one hop at a time based on local link loads and the destination of the demand. Unlike the single-agent approach, no individual agent has access to the full path. Training is conducted centrally, while execution is decentralized.
\end{enumerate}


\subsection{Training Configuration}
\label{sec:setup-hyper}

Table~\ref{tab:hyper} presents the learner configurations. Every optimizer hyperparameter is held equal across the two architectures, so any difference in outcome is attributable to decentralization rather than to tuning. The centralized controller keeps the default settings of Stable-Baselines3 \citep{stable-baselines3}, and MAPPO is adapted to these defaults rather than modifying the baseline to match the proposed method. This approach avoids disadvantaging the baseline. Both methods therefore perform 732 updates over $1.5\times10^{6}$ environment steps using ten random seeds.


\begin{table}[h]
\centering
\small
\caption{Learner configurations for the two architectures, matched on every
optimizer hyperparameter.}
\label{tab:hyper}
\begin{tabular}{lcc}
\hline
& MARL (MAPPO) & Single-agent GNN (PPO) \\
\hline
Environment steps & $1.5\times10^{6}$ & $1.5\times10^{6}$ \\
Rollout buffer & $2048$ & $2048$ ($8\times256$) \\
Policy updates & $732$ & $732$ \\
Learning rate & $3\times10^{-4}$ & $3\times10^{-4}$ \\
Discount $\gamma$ & $0.995$ & $0.995$ \\
GAE $\lambda$ & $0.95$ & $0.95$ \\
Clip range & $0.2$ & $0.2$ \\
Epochs per update & $10$ & $10$ \\
Minibatch & $256$ & $256$ \\
Entropy coefficient & $0.01$ & $0.01$ \\
Value coefficient & $0.5$ & $0.5$ \\
Gradient-norm clip & $0.5$ & $0.5$ \\
Message-passing rounds & $3$ & $3$ \\
Hidden width & $32$ & $64$ \\
Trainable parameters & $17{,}570$ & $268{,}420$ \\
Candidate paths per pair $k$ & --- (hop-by-hop) & $3$ \\
Seeds & $0,\dots,9$ & $0,\dots,9$ \\
\hline
\end{tabular}
\end{table}


The two policies differ in capacity because equalizing parameter counts would require distorting one of the architectures. The decentralized width was selected on training topologies rather than on the evaluation backbones, using held-out TMgen matrices from the seventeen training instances generated with a seed distinct from that used in training.

Both training procedures share identical environment settings: delay penalty $\beta = 0.5$, detour allowance $\sigma = 2$, maximum stretch $\sigma_{\max} = 4$ cost units, reward normalization by the episode's OSPF bottleneck, and a $500$-demand cap on training matrices. The reward function is also consistent, with a flat environment charging a fixed $\beta$ for each detour hop and normalized congestion.

Two key differences remain between the approaches. First, the decentralized policy receives rewards incrementally along the hops of a demand, whereas the centralized policy receives a single reward at the end. As a result, both policies encounter different value targets and effective horizons because the discount factors are the same. Second, the definition of a step is different for both policies. A single centralized step routes an entire demand, while one MARL step advances only a single hop. Therefore, at $1.5\times10^{6}$ environment steps, the centralized controller processes approximately six times as many demands, given that average path length is about six hops. Although the budget is matched in terms of environment steps, it is not matched in terms of routing work. Both differences favor the centralized method, so the decentralized policy's superior performance cannot be attributed to under-training of the baseline.


\subsection{Congestion Regimes}
\label{sec:setup-regimes}
Network simulations are classified into two distinct congestion regimes to account for the differing outcomes produced by the routing problem under each condition. The congestion measure is the bottleneck utilization $U$ of Equations~\eqref{eq:load}--\eqref{eq:bottleneck}, the offered utilization of the busiest link under a given routing, computed as the total demand of the paths traversing a link divided by that link's capacity.

Because it is computed from the demands alone it has no upper bound, and values above $100\%$ express routing pressure that exceeds capacity. The offered load is evaluated using an analytical surrogate to assign regime labels to each traffic matrix. Each test matrix is categorized based on the offered utilization that \emph{OSPF} would produce: \textbf{overload} when $U_{\mathrm{OSPF}}\ge100\%$, and \textbf{feasible} when $U_{\mathrm{OSPF}}<100\%$.

The labelling is confirmed in ns-3: under OSPF the overload matrices lose 14 to 27 percent of packets and the feasible matrices between 0.03 and 0.17 percent.

\subsection{Packet-Level Evaluation Configuration}
\label{sec:setup-ns3}


This study uses ns-3 to evaluate labeled traffic. Each link is a \texttt{PointToPointNetDevice} with the default DropTail transmit queue of 100 packets, and no active queue management, shaping, or congestion control is enabled, so all reported packet loss results solely from buffer overflow. Traffic is one-way UDP, each demand a constant-rate \texttt{OnOff} application with a 1400-byte payload, so loss and delay follow from the routing decision rather than from transport-layer feedback. Propagation delays are installed at nanosecond resolution, which matters on Germany50 where several links fall below 1 ms. Each run spans 8 seconds of simulated time with flows active from 2 to 9 seconds. All utilization figures are adjusted by dividing by 7/8 to account for this. Demand rates and link capacities are both divided by 20, which preserves every utilization, loss, and delay ratio while making the packet count manageable. Last, the complete grid consists of 378 simulations.


\subsection{Performance Metrics}
\label{sec:setup-metrics}

This study evaluates the network performance of the MARL in comparison to other routing methods. Three performance metrics are selected for this evaluation. The first metric is analytical, while the remaining metrics are measured using the ns-3 simulator. All metrics are assessed under both overload and feasible operating regimes.
\begin{enumerate}
  \item \textbf{Maximum offered load.} This metric serves as the training objective and quantifies the extent to which routing exceeds the capacity of the most congested link. It replaces utilization as a metric, since in the ns-3 simulator, traffic cannot exceed 100\% capacity.
  \item \textbf{Packet loss.} This metric represents the proportion of dropped packets, expressed as a percentage. It indicates the impact of congestion on data transmission. Packet loss is captured in percentage.
  \item \textbf{Delay.} This metric measures the mean end-to-end packet transmission time, representing the duration required for a packet to travel from source to destination. Delay is recorded in milliseconds.
\end{enumerate}

% Differences between MARL and OSPF are tested with a two-sided Wilcoxon signed-rank
% test paired on the traffic matrix. Pairing is the appropriate design because each
% learned routing is simulated on exactly the same matrix as the OSPF routing it is
% compared against, so the matrix is held fixed and only the routing varies; matrices
% differ widely in offered load, and an unpaired test would submerge the routing
% effect in that variation. The three model seeds are averaged before pairing, since
% they are three draws of a single policy rather than three independent observations
% of a matrix. Each regime therefore contributes nine paired matrices, three from each
% backbone.
